\chapter{Shared analytic bridges}
\label{chap:shared-analytic-bridges}

The long routes reuse a small collection of shared and branch-specific endpoint interfaces.  This
chapter records them once; it does not make every node a conjunctive premise of every route.  The
displayed repository interfaces are phase-one proposals pending human review, so they are marked
\texttt{notready}; none is Lean-certified here.  In particular, a published input or an upstream
reviewed proof never receives \texttt{leanok} merely because it is cited.

\section{Three presentations of the Poincar\'e constant}

\begin{theorem}[Legacy tests and ambient relaxation have the same finite constants]
  \label{thm:legacy-relaxed-holds-equivalence}
  \notready
  \uses{def:l2-mean-variance,def:smooth-gradient-preenergy,
        def:smooth-gradient-relaxation,def:legacy-smooth-test-poincare}
  Let $\mu$ be a probability law on a finite-dimensional real Hilbert space.  For every finite
  $C\ge0$, the legacy inequality on all globally smooth finite-gradient-energy tests holds at $C$
  if and only if the $L^2$ inequality for the ambient sequential relaxation holds at $C$.
\end{theorem}
\begin{proof}
  Open.  A constant sequence gives relaxed energy at most the raw energy of each realized test, so
  the relaxed inequality implies the legacy one.  Conversely, apply the legacy inequality along
  every approximating sequence.  Variance is continuous under $L^2$ convergence; taking the
  liminf and then the infimum over sequences gives the relaxed inequality.  No equality of raw and
  relaxed energy on an individual test is used.
\end{proof}

\begin{theorem}[Legacy and selected canonical constants agree]
  \label{thm:legacy-canonical-constant-comparison}
  \notready
  \uses{thm:legacy-relaxed-holds-equivalence,
        thm:ambient-intrinsic-relaxation-equality,def:selected-canonical-frontend}
  For every finite-dimensional Borel probability $\mu$,
  \[
    \CP(\mu)=\CPCan(\mu)
  \]
  as extended nonnegative constants, including the value $+\infty$.
\end{theorem}
\begin{proof}
  Open.  The preceding theorem identifies the sets of valid finite constants for the legacy and
  ambient-relaxed presentations.  Their extended optima therefore agree.  The exact
  ambient/intrinsic equality identifies the latter with the selected constant.
\end{proof}

\begin{theorem}[Smooth and locally Lipschitz test constants agree]
  \label{thm:smooth-lipschitz-constant-equivalence}
  \notready
  \uses{def:literature-poincare-constant,def:legacy-smooth-test-poincare}
  If $\mu$ is a finite-dimensional log-concave probability law, then
  \[
    \CP(\mu)=\CPLit(\mu).
  \]
  Gradients and approximation are taken intrinsically on the affine support.
\end{theorem}
\begin{proof}
  Open.  One inequality is restriction to globally smooth tests.  The reverse inequality is the
  smooth-density and energy-recovery theorem on the relative interior of the convex support,
  followed by cutoff and truncation.  This is a Sobolev-density step, not raw/core agreement; see
  \cite{AmbrosioGigliSavare2013Density} for the corresponding weak-gradient framework.
\end{proof}

\begin{theorem}[Selected and literature constants agree]
  \label{thm:selected-literature-constant-equivalence}
  \notready
  \uses{thm:legacy-canonical-constant-comparison,
        thm:smooth-lipschitz-constant-equivalence}
  For every finite-dimensional log-concave probability law $\mu$,
  \[
    \CPCan(\mu)=\CPLit(\mu).
  \]
\end{theorem}
\begin{proof}
  Open.  Both sides equal the legacy smooth-test optimum by the two preceding comparison
  theorems.  The Cheeger-energy identification of
  \cref{thm:smooth-gradient-relaxation-two-cheeger} is not needed.
\end{proof}

\section{Isoperimetric, affine, and one-dimensional endpoints}

\begin{definition}[Isoperimetric Cheeger constant and inverse scale]
  \label{def:isoperimetric-cheeger-scale}
  \notready
  For a probability law $\mu$, let
  \[
    h_\mu=\inf_B
      \frac{\mu^+(B)}{\min\{\mu(B),1-\mu(B)\}},
    \qquad
    \PsiKLS(\mu)=h_\mu^{-1},
  \]
  where the infimum ranges over nontrivial Borel sets and $\mu^+$ is outer Minkowski boundary
  measure in the affine support.  The isoperimetric number $h_\mu$ is not the metric energy
  $\Ch_\mu$ of \cref{thm:smooth-gradient-relaxation-two-cheeger}.
\end{definition}

\begin{theorem}[Cheeger--Poincar\'e comparison for log-concave laws]
  \label{thm:logconcave-cheeger-poincare-comparison}
  \notready
  \uses{def:isoperimetric-cheeger-scale,def:literature-poincare-constant}
  Every finite-dimensional log-concave probability law satisfies
  \[
    \CPLit(\mu)\le4\PsiKLS(\mu)^2,
    \qquad
    \PsiKLS(\mu)^2\le\pi\CPLit(\mu).
  \]
\end{theorem}
\begin{proof}
  This is the published Cheeger and convex Buser--Ledoux comparison, recorded with the
  normalization used here; see \cite{Milman2009Isoperimetric,LeeVempala2018}.  It is an imported
  mathematical input, not a local Lean proof.
\end{proof}

\begin{theorem}[Affine and isotropic KLS formulations]
  \label{thm:affine-isotropic-reduction}
  \notready
  \uses{def:literature-poincare-constant}
  For a fixed $K\ge0$, the following statements are equivalent.
  \begin{enumerate}
    \item Every isotropic log-concave law on every finite-dimensional Euclidean space satisfies
          $\CPLit\le K$.
    \item Every finite-dimensional log-concave law $\mu$ satisfies
          $\CPLit(\mu)\le K\norm{\Cov_\mu}_\op$.
  \end{enumerate}
  In the second statement covariance, whitening, and isotropy are computed on the tangent space of
  the affine support; the zero-dimensional law is included.
\end{theorem}
\begin{proof}
  Open.  The affine statement immediately specializes to covariance identity.  Conversely,
  translate and whiten on the affine support, apply the isotropic statement, and pull the
  inequality back.  The squared operator norm of the inverse whitening map is
  $\norm{\Cov_\mu}_\op$.  This is the standard affine reduction in
  \cite{KannanLovaszSimonovits1995,LeeVempala2018}.
\end{proof}

\begin{theorem}[One-dimensional log-concave Poincar\'e bound]
  \label{thm:one-dimensional-logconcave-poincare}
  \notready
  \uses{def:literature-poincare-constant}
  Every log-concave probability law $\nu$ on a line satisfies
  \[
    \CPLit(\nu)\le4\Var_\nu(T),
  \]
  including the zero-variance case.
\end{theorem}
\begin{proof}
  This is the sharp universal one-dimensional input used by the conditional-fiber route; see
  \cite{Bobkov1999LogConcave}.  It is imported mathematics and has no local \texttt{leanok}.
\end{proof}

These nodes are deliberately modest.  They do not provide a self-adjoint generator, identify a
form gap with the reciprocal of an operator eigenvalue, or prove any of the four long-route gates.
